\chapter{Phase 2 Plan}
\label{ch:phase2}

Phase~1 delivered the camera-node vision core. Phase~2 builds the central
system that turns confirmed plates into decisions and moves the node onto
embedded hardware. The work packages below follow the objectives of
Chapter~\ref{ch:intro}.

\section{Work Packages}

\subsection{Vehicle verification}

A central database will hold, for each registration, the registration status,
insurance, pollution certificate, stolen and blacklist status, and previous
sightings. Each recognised plate is looked up and the vehicle classified as
\emph{authorised}, \emph{non-compliant} (for example an expired document),
\emph{unauthorised}, \emph{blacklisted} or \emph{stolen}. Because the reader can
make a one-character error, the lookup will also return near matches with the
reading confidence so that an operator can confirm doubtful cases.

\subsection{Alerts and dispatch}

A real-time dashboard will show the vehicle details, the captured image and the
timestamp, with an administrative interface for managing the database and
users. Stolen and blacklisted vehicles are forwarded to a dispatch module that
finds the nearest available response officer by road-network distance.

\subsection{Multi-node correlation}

Each sighting is stored with its camera node and time. Queries over the stored
sightings will detect convoy movement (vehicles seen together at several
nodes), route anomalies (a vehicle at a node it should not reach in the time
available) and loitering, and will answer retroactive queries for the movement
history of a plate.

\subsection{Embedded deployment}

Detection, restoration and recognition will be ported to the Jetson Nano-class
node. The tests are: end-to-end alert latency from the vehicle passing to the
alert on the dashboard; memory and speed of the UVMNet-based restorer on the
chosen module (the Mamba kernel requires an NVIDIA GPU); and field tests in
clear weather, rain and fog, including night operation with the NoIR camera and
IR illuminator.

\subsection{Completing the weather stage}

Three items of Phase~1 remain open: the effect of deraining on plate detection
and reading in real rain (comparing detection rate, character-level accuracy and
complete-plate rate on the same samples before and after restoration), routing
of frames to the restoration models by a weather classifier in the live loop,
and more labelled footage from fog, rain and night, since labelled real
footage was the data that moved the reader most.

\section{Evaluation Plan}

Table~\ref{tab:eval} lists the measures with which the complete system will be
evaluated.

\begin{table}[t]
\centering
\caption{Planned evaluation of the complete system.}
\label{tab:eval}
\small
\begin{tabularx}{\textwidth}{@{}lXl@{}}
\toprule
\textbf{Measure} & \textbf{Definition} & \textbf{Condition}\\
\midrule
Plate recognition rate & Plates read exactly / plates passing & clear, fog, rain, night\\
Restoration benefit & Change in recognition rate with and without restoration & fog, rain\\
Correlation accuracy & Correct convoy / anomaly / loitering decisions & multi-node\\
Alert latency & Time from the vehicle passing to the alert shown & node and network\\
Node load & Frames per second, memory and power on the node & Jetson module\\
\bottomrule
\end{tabularx}
\end{table}

\section{Risks}

\begin{itemize}
  \item \textbf{Reading accuracy.} Small, blurred plates and non-Kerala series
        are read less reliably; the database lookup and operator confirmation
        must tolerate single-character errors.
  \item \textbf{Embedded speed.} UVMNet is large (116.2 million parameters)
        and needs the CUDA Mamba kernel; if it is too slow on the node, the
        restoration can run only on frames flagged as foggy, or on a central
        GPU.
  \item \textbf{Data.} Public Indian plate data are exhausted for our purpose;
        further gains need self-collected labelled footage.
\end{itemize}
